\Needspace{16\baselineskip}
\section{Proof of Theorem~\ref*{thm:sharper}(2): coverage of real thresholds}\label{sec:thresholdproof}

This appendix proves the uniform threshold statement in
\cref{thm:sharper}(2). We construct graph pairs covering overlapping
ranges of separating thresholds, then extend this coverage to every
real threshold from $6.0101$ by joining complete graphs.

This appendix combines an analytical construction with computer-assisted
checks of finite graph data and threshold inequalities. OpenAI Codex
assisted with exploring the constructions and with developing and
reviewing the verification code. The finite threshold comparisons and
interval overlaps are checked using exact integer and rational
arithmetic; the extension to all thresholds in the claimed range is
proved analytically.

\begin{theorem}[Failure at every real threshold from $6.0101$]
\label{thm:real-thresholds}
For every real $t\ge6.0101$, there are finite simple graphs $G,H$ such that
\[
 \boxed{\C(G\times H)<t<\min\{\C(G),\C(H)\}.}
\]
In particular, $f_C(t)<t$ throughout this range.
\end{theorem}


The proof gives the stronger range $t>T_0$, where
$T_0=6.0100886404\ldots$ is defined exactly
in~\eqref{eq:uniform-optimized}. We first optimize a counterexample
near six, then use eight witnesses to cover a full unit interval and
extend that cover by joining complete graphs. The optimization is
analytical, and all comparisons for the stated decimal cutoff use
exact integer arithmetic.

We use the weighted product bound~\eqref{eq:weighted-main} from
\suppref{sec:interpolation}.
Throughout this section, put
\[
 A=U_6,\qquad D=\overline A,\qquad
 h(X)=\phif{\F_2}(X),\qquad q(X)=\vartheta(\overline X).
\]
Both $h$ and $q$ are OR spectral points. The constructions in
\cref{lem:binary,lem:clebsch} give
\[
 h(A)=6,\qquad q(D)\le4,\qquad
 \C(Q)\ge16^{1/3},\qquad q(Q)\le8/3.
\]
We first provide the remaining small-graph estimates.

\subsection{Fractional colourings and a 37-vertex partner}

We use the fractional-colouring definitions in
\hyperref[def:fractional-colouring]{Supplementary App.~\ref*{sec:spectrum}}
and the inequality $h(X)\le\chi_f(X)$ proved in
\eqref{eq:haemers-colouring} by the coordinate-vector fitting construction.

In particular,
\begin{equation}\label{eq:threshold-colouring-D}
 h(D)\le\chi_f(D)\le\frac{16}{3}.
\end{equation}
Indeed, $A$ is vertex-transitive. Coordinate permutations and the maps
$T_v(x)=x+(x\cdot v)v$ for even-weight $v$ preserve dot products
and odd weight; moreover, $T_v^2$ is the identity. Taking
$v=e_1+e_2+e_3+e_4$ or $v=\mathbf1$ sends $e_1$ to a vector of
weight three or five. Permutations then reach every odd-weight vector.
Thus $D$ is also vertex-transitive. Averaging its independent six-set
$\{e_1,\ldots,e_6\}$ over the automorphism group covers every vertex
equally. Normalizing this coverage to one gives a fractional colouring
of total weight $32/6$, as explained in \suppref{sec:spectrum}, and proves
\eqref{eq:threshold-colouring-D}.

\begin{lemma}[A self-complementary circulant]\label{lem:circulant37}
Let $R$ be the graph on $\mathbb Z/37\mathbb Z$ in which two vertices
are adjacent when their undirected cyclic distance belongs to
\[
 S=\{6,7,8,9,10,12,13,15,18\}.
\]
Then
\[
 \C(R)=q(R)=\sqrt{37},\qquad h(R)\le\chi_f(R)=\frac{37}{6}.
\]
\end{lemma}
\begin{proof}
Multiplication by six exchanges the two members of each distance pair
\[
 (1,6),(2,12),(3,18),(4,13),(5,7),
 (8,11),(9,17),(10,14),(15,16).
\]
Since $S$ chooses one member of every pair, this map is an isomorphism
$R\cong\overline R$. Consequently the 37 words $(i,6i)$ form a clique
in $R\OR R$, proving $\C(R)\ge\sqrt{37}$. The graph is
vertex-transitive, so Lov\'asz's identity
$\vartheta(R)\vartheta(\overline R)=37$~\cite{Lovasz}, together with
self-complementarity, gives the matching theta upper bound.

The set $B=\{0,6,12,18,24,30\}$ is a six-clique because its nonzero
undirected differences lie in $\{6,12,18,13,7\}$. Thus $6B$ is an
independent six-set. Its 37 translates cover every vertex six times,
giving $\chi_f(R)\le37/6$. Conversely, every independent set has
size at most $\vartheta(R)=\sqrt{37}<7$, hence at most six. Summing
the covering constraints of any fractional colouring gives
$37\le6\chi_f(R)$. Thus $\chi_f(R)=37/6$, proving the claim.
\end{proof}

\subsection{The initial pair near six}

Put $G=A\sqcup(D\join K_1)$. For any $\varphi\in\spec$, write
$x=\varphi(A)$ and $y=\varphi(D)$. The 32 diagonal words in
$A\OR D$ form a clique, so $xy\ge32$. The join and maximum rules give
$\varphi(G)=\max\{x,y+1\}$. If this maximum is $r$, then
$32\le xy\le r(r-1)$, and hence
\begin{equation}\label{eq:near-six-factors}
 \C(G)\ge\frac{1+\sqrt{129}}2>\sqrt{37}=\C(R).
\end{equation}
The projections give the spectral upper bounds
\[
\begin{array}{c|cc}
 &h&q\\ \hline
 A\times R&6&\sqrt{37}\\
 (D\join K_1)\times R&37/6&5.
\end{array}
\]
The geometric bound in~\eqref{eq:weighted-main} gives
\[
 \C(G\times R)\le
 \min_{0\le\lambda\le1}
 \max\{6^\lambda(\sqrt{37})^{1-\lambda},
          (37/6)^\lambda5^{1-\lambda}\}.
\]
Put $d=\log(\sqrt{37}/6)>0$ and $\ell=\log(6/5)>0$,
using natural logarithms.
After subtracting $\log6$, the logarithms of the two branches are
$(1-\lambda)d$ and $2\lambda d-(1-\lambda)\ell$.
The first decreases and the second increases with $\lambda$.
They agree at $\lambda_*=(\ell+d)/(\ell+3d)\in(0,1)$,
which therefore minimizes their maximum. Consequently
\begin{equation}\label{eq:uniform-optimized}
 \C(G\times R)\le T_0
 :=6\exp\!\left(\frac{2d^2}{\ell+3d}\right)
 =6.0100886404\ldots .
\end{equation}
Together with~\eqref{eq:near-six-factors}, this proves failure at every
threshold $T_0<t<\sqrt{37}$.

To prove $T_0<6.0101$ exactly, without using the decimal approximation
in~\eqref{eq:uniform-optimized}, use $\lambda=93/106$ in the geometric
bound. Direct integer exponentiation gives
\begin{equation}\label{eq:uniform-integer-check}
\begin{aligned}
 6^{186}37^{13}10000^{212}&<60101^{212},\\
 37^{93}5^{13}10000^{106}&<6^{93}60101^{106}.
\end{aligned}
\end{equation}
Taking the positive $212$th and $106$th roots, respectively, shows that
both branches are strictly below $60101/10000=6.0101$.
The supplementary script \nolinkurl{verification/check_uniform_cutoff.py}
in the accompanying repository~\cite{ZeissData} checks exactly these
two integer comparisons. The same repository includes
\nolinkurl{verification/check_exact.py} and
\nolinkurl{verification/audit_uniform_independent.py}, which check the
finite graph inputs and interval comparisons used below.

For the interval cover below, we retain a simpler arithmetic bound.
Use \eqref{eq:weighted-main} with $\lambda=13/15$ and
$\sqrt{37}<73/12$. The two arithmetic means agree, giving
\[
 \C(G\times R)
 \le\frac{13}{15}\,6+\frac{2}{15}\frac{73}{12}
 =\frac{13}{15}\frac{37}{6}+\frac{2}{15}\,5
 =\frac{541}{90}.
\]
Thus this pair also supplies the interval $(541/90,\sqrt{37})$.
The next three subsections cover every $t>541/90$; combining that
cover with the optimized initial interval proves the theorem.

\subsection{The cycle inputs}

Write $C_n$ for the cycle on $\mathbb Z/n\mathbb Z$, with two vertices
adjacent when their difference is $1$ or $-1$ modulo $n$ ($n\ge3$).
For the pentagon $C_5$, the standard theta calculation and its five-word
square code give
\[
 \C(C_5)=q(C_5)=\sqrt5,\qquad h(C_5)\le\chi_f(C_5)=5/2.
\]
Explicitly, the words $(i,2i)$, for $i\in\mathbb Z/5\mathbb Z$, form
a clique in $C_5\OR C_5$. The upper bound is Lov\'asz's pentagon
calculation~\cite{Lovasz}; the five translates of an independent pair
give the fractional colouring.

Write $Z=\overline{C_9}$. We need only
\[
 \C(Z)\ge81^{1/3},\qquad h(Z)\le9/2,\qquad q(Z)<22/5.
\]
For the first inequality, take the 81 words
\[
 (x,y,2x+4y),\qquad x,y\in\mathbb Z/9\mathbb Z.
\]
If two distinct words were nonadjacent in $Z^{\OR3}$, their coordinate
differences would have representatives $a,b,c\in\{-1,0,1\}$ with
$c-2a-4b\equiv0\pmod9$. Its absolute value is at most seven, so this
is an integer equality. Since $4>2+1$ and $2>1$, it forces
$a=b=c=0$, a contradiction. The nine translates of an edge of $C_9$
are independent pairs of $Z$, proving the colouring bound $9/2$.
Finally, the odd-cycle theta formula~\cite{Lovasz} gives
\[
 q(Z)=\vartheta(C_9)
 =\frac{9\cos(\pi/9)}{1+\cos(\pi/9)}<\frac{22}{5}.
\]
For a rational verification of the strict inequality, put $x=\pi/9$.
Since $1/3<x<1$, the alternating Taylor estimate gives
\[
 \cos x\le1-\frac{x^2}{2}+\frac{x^4}{24}
 \le1-\frac{11x^2}{24}<\frac{205}{216},\qquad
 \frac{9(205/216)}{1+205/216}=\frac{1845}{421}<\frac{22}{5}.
\]

\subsection{Balancing the bounds and adding cliques}

For an auxiliary graph $J$, define
\[
 G_J=A\sqcup(D\join J),\qquad
 \rho_b=\frac{b+\sqrt{b^2+128}}2.
\]
If $\C(J)\ge b>0$, then \cref{lem:joined-complement} gives
$\C(G_J)\ge\rho_b$. The following lemma chooses the interpolation
weight once and for all from three projection bounds.

\begin{lemma}[Balancing and adding cliques]\label{lem:threshold-padding}
Let $\C(J)\ge b>0$, and let $H$ be a partner graph. Suppose
$a>6$ and $s>d$ satisfy
\[
 s\ge q(H),\qquad
 a\ge\min\{h(D\join J),h(H)\},\qquad
 d\ge q(D\join J).
\]
Put
\begin{equation}\label{eq:threshold-balanced}
 U=\frac{as-6d}{a-6+s-d}.
\end{equation}
For every integer $j\ge0$, there are graphs $G^{(j)},H^{(j)}$ with
\[
 \C(G^{(j)}\times H^{(j)})\le U+j,\qquad
 \C(G^{(j)})\ge\rho_b+j,\qquad
 \C(H^{(j)})=\C(H)+j.
\]
\end{lemma}
\begin{proof}
Take
\[
 G^{(j)}=(A\join K_j)\sqcup(D\join J\join K_j),\qquad
 H^{(j)}=H\join K_j,
\]
omitting $K_0$. The categorical product distributes over disjoint
union. Projection to the factors and additivity under join give the
following bounds on its two components.
\[
\begin{array}{c|cc}
 &h&q\\ \hline
 (A\join K_j)\times H^{(j)}&6+j&s+j\\
 (D\join J\join K_j)\times H^{(j)}&a+j&d+j.
\end{array}
\]
Choose $\lambda=(s-d)/(a-6+s-d)\in(0,1)$. By construction,
\[
 6\lambda+s(1-\lambda)=a\lambda+d(1-\lambda)=U.
\]
Apply the spectral interpolation inequality~\eqref{eq:interpolate}
to $h,q$ with these weights. The resulting single spectral point,
the maximum rule, and the weighted
arithmetic--geometric mean bound the product capacity by $U+j$.
Finally, every spectral point takes value $\varphi(G_J)+j$ on
$G^{(j)}$ and value $\varphi(H)+j$ on $H^{(j)}$.
Taking spectral minima gives the two factor claims.
\end{proof}

\subsection{A single coverage table}

In \cref{tab:threshold-cover}, $kC_5$ means the join of $k$ copies of $C_5$,
and every terminating decimal denotes an exact rational number.
The columns $s,a,d$ satisfy the preceding lemma, and $U$ is
given by~\eqref{eq:threshold-balanced}. Each row gives
\begin{equation}\label{eq:threshold-row}
 \C(G_J\times H)\le U<L<\min\{\C(G_J),\C(H)\}.
\end{equation}
Thus that row works for every threshold in the open interval $(U,L)$.

\begin{table}[htbp]
\centering
\small
\renewcommand{\arraystretch}{1.18}
\begin{tabular}{@{}llccccc@{}}
\toprule
$J$ & $H$ & $s$ & $a$ & $d$ & $U$ & $L$\\
\midrule
$K_1$ & $R$ & $73/12$ & $37/6$ & $5$ & $541/90$ & $6.08$\\
$K_1$ & $K_4\join C_5$ & $156/25$ & $19/3$ & $5$ & $357/59$ & $6.17$\\
$K_2$ & $K_4\join C_5$ & $156/25$ & $13/2$ & $6$ & $228/37$ & $6.23$\\
$K_2$ & $K_2\join Z$ & $32/5$ & $13/2$ & $6$ & $56/9$ & $6.325$\\
$K_2$ & $K_2\join2C_5$ & $162/25$ & $7$ & $6$ & $234/37$ & $6.47$\\
$K_2$ & $3C_5$ & $168/25$ & $22/3$ & $6$ & $498/77$ & $6.7$\\
$C_5$ & $K_7$ & $7$ & $7$ & $156/25$ & $289/44$ & $6.88$\\
$Q$ & $K_1\join R$ & $85/12$ & $43/6$ & $20/3$ & $775/114$ & $7.02$\\
\bottomrule
\end{tabular}
\caption{One formula certifies every interval $(U,L)$. The first factor
is $G_J=A\sqcup(D\join J)$; $s,a,d$ are spectral projection bounds.}
\label{tab:threshold-cover}
\end{table}

Here $s$ bounds $q(H)$ and $d$ bounds $q(D)+q(J)$.
For $a$, use $h(D)+h(J)$ in rows 2 and 6, using
\eqref{eq:threshold-colouring-D}, and $h(H)$ in all other rows.
All entries follow from the preceding graph bounds, additivity under
join, and the elementary enclosures
\[
 2.235<\sqrt5<2.24,\qquad
 2.5<16^{1/3},\qquad 4.325<81^{1/3},
\]
and $6.08<\sqrt{37}<73/12$. Each follows by squaring or cubing.

For the first factor, $L<\rho_b$ follows from $L>b$ and
$L(L-b)<32$. For each choice of $J$, it suffices to check the largest
$L$. Using the respective lower bounds $b=1,2,2.235,2.5$, we obtain
\[
\begin{aligned}
 6.17(6.17-1)&=31.8989<32,&
 6.7(6.7-2)&=31.49<32,\\
 6.88(6.88-2.235)&=31.9576<32,&
 7.02(7.02-2.5)&=31.7304<32.
\end{aligned}
\]
For the partners, the capacity lower bounds are, in row order,
\[
 \sqrt{37},\ 4+\sqrt5,\ 4+\sqrt5,\ 2+81^{1/3},\
 2+2\sqrt5,\ 3\sqrt5,\ 7,\ 1+\sqrt{37}.
\]
The same rational enclosures place each strictly above its listed $L$.
This proves~\eqref{eq:threshold-row} for every row.

\FloatBarrier
\begin{proof}[Completion of the proof of \cref{thm:real-thresholds}]
Write $(U_i,L_i)$ for the intervals in
\cref{tab:threshold-cover}. Direct cross-multiplication gives
$U_i<L_i$ and $U_{i+1}<L_i$; the smallest consecutive overlap
is $6.325-234/37=1/1480>0$. Their union therefore contains
$(541/90,7.02)$, an interval of length greater than one.

For any $t>c:=541/90$, choose $j=\lceil t-c\rceil-1\ge0$.
Then $c<t-j\le c+1<7.02$, so some row has $U_i<t-j<L_i$.
\Cref{lem:threshold-padding} gives
\[
 \C(G^{(j)}\times H^{(j)})\le U_i+j<t<L_i+j
 <\min\{\C(G^{(j)}),\C(H^{(j)})\}.
\]
Thus the same eight witnesses, extended by joining complete graphs,
cover every $t>541/90$.

The optimized initial pair covers $(T_0,\sqrt{37})$, and
\eqref{eq:uniform-integer-check} gives
\[
 T_0<6.0101<541/90<\sqrt{37}.
\]
The two ranges overlap, so together they cover every $t>T_0$.
In particular they cover every $t\ge6.0101$, including the stated
endpoint. This proves the uniform clause of \cref{thm:sharper}.
\end{proof}

For the real-threshold spectral claim in \cref{sec:spectral-consequences}, fix
$t\ge6.0101$ and choose a pair $G_t,H_t$ from the construction above.
By~\eqref{eq:duality}, there is a spectral point $\psi_t\in\spec$
attaining the capacity of $G_t\times H_t$. Since every spectral point
bounds capacity from above, we obtain
\[
 \psi_t(G_t\times H_t)=\C(G_t\times H_t)<t
 <\min\{\C(G_t),\C(H_t)\}
 \le\min\{\psi_t(G_t),\psi_t(H_t)\}.
\]
This is Simonyi's implication~\cite[Thm.~8]{Simonyi} applied to each
pair. The choice of $\psi_t$ may depend on $t$.

\subsection{Scope of the sharper thresholds}

The positive cases $1$ and $2$ recalled in \suppref{sec:spectrum},
together with \cref{thm:finite}, show that the smallest integer value of
the smaller factor capacity admitting a counterexample lies between
$3$ and $7$. The unresolved exact integer minima are $3,\ldots,6$. The pair in
\cref{sec:subsixproof} has smaller factor capacity in
$[\sqrt{27},5.1965612]$; its exact value is unknown.
The earlier pair in~\eqref{eq:near-six-factors} has exact smaller
factor capacity $\sqrt{37}$.
By \cref{subsix:thm:main}, the sub-six pair gives failure at every
threshold $5.193957\le t<\sqrt{27}$, whereas this section covers
every $t>T_0$. A cover joining the lower interval to this upper ray
is still missing. In particular, the strict-threshold formulation
does not include $t=\sqrt{27}$ merely from the lower bound
$\min\{\C(G),\C(H)\}\ge\sqrt{27}$, and the upper-ray proof does not
include $t=T_0$. This does not assert that each individual threshold
between these endpoints is unresolved.

Neither result classifies all possible exact noninteger minimum
capacities; no optimality is claimed.
Polak's circular-graph constructions provide, after complementation,
graphs whose OR-capacities approach any integer $r\ge3$ from
below~\cite[Chap.~9, Thms.~9.3.1 and~9.3.3]{PolakThesis}.
These constructions do not determine whether the product identity holds
at those capacities.
